
	\newpage
	\section{Lineare Gleichungssysteme} \label{chap:LGS}


	\defi{ \n
		Mit linearen Gleichungssystemen muss man sich des Öfteren auseinandersetzen, da sie alleine schon als Teilproblem häufig auftreten.
		Wann immer wir ein lineares Gleichungssystem (kurz LGS) haben, berechneten wir dessen Lösung bisher mithilfe des \fat{Gauß-Elminationsverfahrens}.
		\n
		Was war doch gleich ein LGS?
		\n
		Ein LGS haben wir immer dann, wenn wir mehrere Gleichungen mit Unbekannten haben und wissen möchten, welche Werte die einzelnen Unbekannten haben müssen/können, damit alle Gleichungen erfüllt sind
		\n 
		Allgemein beschreiben wir sie in der Form $Ax=b$, wobei $A$ eine Matrix und $x$ und $b$ Vektoren sind ($b$ wird auch manchmal "Ergebnisvektor" genannt).
	}
	
	\sbreak
	\bsp{
		\begin{align*}
			1x + 2y + 2z &= 3 &&(1)\\
			2x + 1y - 2z &= 2 &&(2)\\
			3x + 0y + 2z &= 6 &&(3)
		\end{align*}
		Zu lösen: \\Welche Werte müssen $x$, $y$ und $z$ haben, damit alle Gleichungen $(1)$ -- $(3)$ erfüllt sind?
	}
	
	\lbreak
	% =========================================================	
	\subsection{Gauß-Elimination} \label{chap:gausselim}
	% =========================================================
	
	\methode{
		\n
		Wir führen alle Gleichungen in eine Matrix-Schreibweise um und führen eine \fat{Vorwärtssubstitution} aus (das nennt man den Vorgang, die Matrix in eine rechts-obere Dreiecksmatrix zu bringen).
		\n
		In dieser Vorwärtssubstitution verwenden wir die typischen drei elementaren Zeilenoperationen:
		\begin{enumerate}
			\item Reihe + Reihe
			\item Reihe multipliziert mit Skalar
			\item Vertauschen von 2 Reihen
		\end{enumerate}
	}
	
	\newpage
	\bsp { (die Zahlen aus vorherigem Beispiel)
		\begin{align*}
			\begin{bmatrix}
				1 & 2 & 2 & |&3 \\
				2 & 1 & -2 & |&2 \\
				3 & 0 & 2 & |&6
			\end{bmatrix} = 
			\begin{bmatrix}
				1 & 2 & 2 & |&3\\
				0 & -3 & -6 & |&-4\\
				0 & -6 & -4 & |&-3
			\end{bmatrix} = 
			\begin{bmatrix}
				1 & 2 & 2 & |&3\\
				0 & -3 & -6 & |&-4\\
				0 & 0 & 8 & |&5
			\end{bmatrix}
		\end{align*}
		Ab diesem Zeitpunkt verwenden wir jetzt nur noch die \fat{Rückwärtssubstitution} um auf alle Variablenwerte zu schließen, indem wir uns von unten nach oben vorarbeiten.
		\begin{align*}
			8 \cdot z &= 5 \\
			z &= \frac{5}{8} \\
			&\\
			-3 y - 6 z &= -4 \\
			-3 y - 6 \cdot \frac{5}{8} &= -4 \\
			-3 y &= -4 + \frac{30}{8} \\
			-3 y &= \frac{-1}{4} \\
			y &= \frac{1}{12} \\
			& \\
			1 x + 2 y + 2 z &= 3 \\
			1 x + 2 \frac{1}{12} + 2 \frac{5}{8} &= 3 \\
			x &= 3 - \frac{1}{6} - \frac{10}{8} \\
			x &= \frac{19}{12}
		\end{align*}
		%$z = \nicefrac{5}{8}$, dann in zweite Zeile einsetzen, $y$ ausrechnen, etc\dots
		Wir konnten die (in diesem Fall drei) Unbekannten $x$, $y$ und $z$ so bestimmen, dass sie alle Gleichungen erfüllen.
	}
	
	\newpage
	\subsection{Pivotisierung}
	
	\defi{ %\fat{Pivotisierung}
		\n
		Jetzt möchten wir dieses Verfahren einem Computer beibringen. Als \fat{Laufzeit} erhalten wir für Gauß (= Vorwärtssubst.) und Rückwärtssubst. zusammen $\bigO(n^3)$. %$\bigO(n^2)$ und für die Rückwärtssubst. $\bigO(n)$.
		\n
		Wir wissen schon aus vergangenen Tutorübungen, dass Zahlen von einem PC häufig nicht korrekt dargestellt werden (können) und deswegen gerundet werden müssen (siehe \refChapter{chap:rundung}).
		\n
		\noindent Was dagegen tun?
		\n
		Als Pivotelement wird das Element \fat{ganz links oben} in dem noch übrig gebliebenen Teil der Matrix genannt.
		Ein betragsmäßig kleines Pivot führt zu einem instabileren Algorithmus, da die Faktoren beim Verrechnen von zwei Zeilen dann größer werden.
		\n
		Somit ist es immer besser, die Zeilen der Matrix so zu vertauschen, dass das betragsmäßig größte Element zum Pivot wird.
		Wenn wir dies bei jeder Iteration (Spalte) machen, erhalten wir die maximale Stabilität.
	}

	\sbreak
	\wichtig{ Kurze Überlegung: \n
		 Wenn das Pivot 0 ist, keine Pivotisierung angewendet wurde und mindestens ein anderer Wert in der ersten Spalte ungleich 0 ist, kann der Algorithmus gar nicht weiterarbeiten! \n
	 	Wer mir nicht glaubt, probiere mal, die zweite Zeile mit der ersten so zu verrechnen, dass eine neue 0 entsteht \Laughey.
 	}
